\section*{Capítulo \ref{ch:b3:relativistic-dynamics} --- Dinámica relativista}

\begin{solution}{exo:b3:relativistic-dynamics:1}
Electrón: $\gamma = 7.09$; $pc = \gamma\beta\,mc^2 = \qty{3.59}{MeV}$,
luego $p = \qty{3.59}{MeV}/c$; $E = \qty{3.62}{MeV}$, $E_k =
\qty{3.11}{MeV}$. Protón: $E = 938 + 1000 = \qty{1938}{MeV}$, $\gamma
= 2.07$, $pc = \sqrt{E^2 - (mc^2)^2} = \qty{1696}{MeV}$, $\beta =
pc/E = 0.875$.
\end{solution}

\begin{solution}{exo:b3:relativistic-dynamics:2}
(a) $\num{e-3} \times \num{9e16} = \qty{9e13}{J}$. (b) $\num{6e13}/
\num{9e16} \approx \qty{0.7}{g}$ --- la bomba convirtió menos de un
gramo. (c) $L/c^2 = \qty{4.2e9}{kg/s}$; a lo largo de \qty{e10}{yr}
($\num{3.2e17}{}\,\unit{s}$): $\num{1.3e27}{}\,\unit{kg}$, alrededor del
$0.07\%$ del Sol. (d) $\num{5e4}/\num{9e16} \approx
\qty{0.6}{ng}$ --- real y para siempre impesable.
\end{solution}

\begin{solution}{exo:b3:relativistic-dynamics:3}
(a) $p = P/c$ por segundo: una fuerza de retroceso $F = P/c = \qty{1.7e-11}{N}$.
(b) La reflexión duplica la transferencia: $F = 2\Phi A/c =
2 \times 1400 \times 100/\num{3e8} \approx \qty{0.9}{mN}$. (c) $a
\approx \qty{9e-5}{m/s^2}$; al cabo de \qty{2.6e6}{s}: $v \approx
\qty{240}{m/s}$ --- lenta al arrancar, pero el motor nunca se queda seco. (d)
El momento de la luz solar (junto con el viento solar) empuja el polvo
continuamente hacia fuera: la cola se aleja del Sol, no queda detrás del
cometa.
\end{solution}

\begin{solution}{exo:b3:relativistic-dynamics:4}
(a) $m = E_0/c^2$: $\num{0.511e6} \times \num{1.6e-19}/\num{9e16} =
\qty{9.1e-31}{kg}$. (b) $\num{e9} \times \num{1.6e-19}/\num{3e8} =
\qty{5.3e-19}{kg.m/s}$. (c) $\sqrt s = \qty{4}{GeV}$ de energía en el
centro de momentos. (d) Con $c = 1$, la energía, el momento y la
masa comparten una sola unidad y toda fórmula pierde sus $c$; para volver al
SI, se reinserta la única potencia de $c$ que ajusta las dimensiones.
\end{solution}

\begin{solution}{exo:b3:relativistic-dynamics:5}
(a) $2E_k/mc^2$: $1.96$, $3.9$, $17.6$, $59$ --- absurdo a partir del
primero. (b) $1 - (1 + E_k/mc^2)^{-2}$: $0.74$, $0.89$, $0.990$,
$0.9989$. (c) $8.4/0.99946c = \qty{28.0}{ns}$; los $7.7c$ de Newton
habrían dado \qty{3.7}{ns}. (d) El calorímetro demostró que los electrones
llevaban realmente los \qty{15}{MeV} enteros hasta el blanco: la energía
estaba toda allí y, aun así, la velocidad había dejado de crecer --- el
trabajo compra ahora momento y energía, no velocidad.
\end{solution}

\begin{solution}{exo:b3:relativistic-dynamics:6}
(a) $m_\pi c^2 = \sqrt{p^2c^2 + m_\mu^2c^4} + pc$; aíslese la raíz
y elévese al cuadrado: $pc = (m_\pi^2 - m_\mu^2)c^4/2m_\pi c^2 =
(139.6^2 - 105.7^2)/(2 \times 139.6) = \qty{29.8}{MeV}$. (b) $E_\mu =
\sqrt{29.8^2 + 105.7^2} = \qty{109.8}{MeV}$: $E_k = \qty{4.1}{MeV}$,
$\beta = 29.8/109.8 = 0.27$. (c) $E_\nu = pc = \qty{29.8}{MeV}$. (d)
Al pasar al sistema del laboratorio: $E_{\text{lab}} = \gamma(E^* \pm \beta p^*c)$ con $\beta
\approx 1$: entre $50 \times (109.8 - 29.8) = \qty{4.0}{GeV}$ y
$50 \times (109.8 + 29.8) = \qty{7.0}{GeV}$.
\end{solution}

\begin{solution}{exo:b3:relativistic-dynamics:7}
(a) Para dos cuantos sin masa, $M^2c^4 = 2E_1E_2(1 - \cos\theta) =
2 \times 200^2 \times \tfrac12$: $M = \qty{200}{MeV}/c^2$. (b)
Exactamente $m_{\pi^0}$ --- la \omterm{def:b3:relativistic-dynamics:four-momentum}{masa invariante} es la masa de la partícula,
en todos los sistemas. (c) Se calcula $M$ para cada pareja de fotones del
suceso: las parejas al azar se reparten suavemente y las hijas verdaderas de
una partícula se amontonan en su masa --- la joroba. (d) $\theta = 0$:
$M = 0$; un haz de luz paralelo no tiene sistema propio --- se mueve a $c$
como un todo, igual que cada uno de sus fotones.
\end{solution}

\begin{solution}{exo:b3:relativistic-dynamics:8}
(a) $\gamma = \num{6.8e12}/\num{9.38e8} = 7250$; $c - v \approx
c/2\gamma^2 = \qty{2.9}{m/s}$. (b) $f = v/L \approx \num{3e8}/27000 =
\qty{11.1}{kHz}$: once mil vueltas por segundo. (c) $\num{3e14}
\times \qty{6.8}{TeV} = \qty{3.3e8}{J} \approx \qty{78}{kg}$ de TNT
--- en un haz más fino que un cabello. (d) $2 \times \num{3.3e8}/\num{9e16} \approx
\qty{7}{\micro g}$: las existencias de materia por nacer.
\end{solution}

\begin{solution}{exo:b3:relativistic-dynamics:9}
(a) La \omterm{def:b3:relativistic-dynamics:four-momentum}{masa invariante} de un fotón es $0$; la de un par $e^+e^-$ vale al
menos $2m_{\text{e}}c^2$. El invariante no puede cambiar en una
desintegración aislada: prohibida. (Equivalentemente: en ningún sistema
puede equilibrarse el momento). (b)
Frontalmente, $M^2c^4 = 4E\epsilon$; la creación de pares exige $Mc^2 \ge
2m_{\text{e}}c^2$, es decir, $E\epsilon \ge (m_{\text{e}}c^2)^2$. (c)
Otro fotón de \qty{511}{keV}; para \qty{50}{keV}, un compañero de
$0.511^2/0.05 = \qty{5.2}{MeV}$. (d) Un fotón de TeV encuentra el umbral
$(m_{\text{e}}c^2)^2/E \sim \qty{0.3}{eV}$ en la luz estelar ordinaria: el
propio cielo absorbe los rayos $\gamma$ de TeV a distancias cosmológicas
--- el universo no es transparente a todas las energías.
\end{solution}

\begin{solution}{exo:b3:relativistic-dynamics:10}
(a) $P_{e'} = P_\gamma + P_e - P_{\gamma'}$; elévense al cuadrado ambos
miembros (invariantes): $m_{\text{e}}^2c^4 = m_{\text{e}}^2c^4 + 2P_\gamma\!
\cdot\!P_e - 2P_{\gamma'}\!\cdot\!P_e - 2P_\gamma\!\cdot\!P_{\gamma'}$.
(b) Con $P_\gamma\!\cdot\!P_e = E m_{\text{e}}$,
$P_{\gamma'}\!\cdot\!P_e = E'm_{\text{e}}$ y $P_\gamma\!\cdot\!
P_{\gamma'} = EE'(1 - \cos\theta)/c^2$:
$m_{\text{e}}c^2(E - E') = EE'(1 - \cos\theta)$, que en longitudes de onda
($E = hc/\lambda$) es el desplazamiento enunciado. (c) $h/m_{\text{e}}c =
\qty{2.43}{pm}$, y como mucho \qty{4.9}{pm}: una parte en $10^5$ de la luz
visible (invisible), y un porcentaje de un rayo X de \qty{100}{pm}
(medible). (d) Que un fotón choca como una bola de billar ---
transportando un momento $h/\lambda$ intercambiado en sucesos
\emph{individuales}, y no solo energía a granel.
\end{solution}

\begin{solution}{exo:b3:relativistic-dynamics:11}
(a) $M^2c^4 = m_{\text{p}}^2c^4 + 4E\epsilon$ frontalmente (con el protón
ultrarrelativista); umbral en $M = m_\Delta$. (b) $E =
(1232^2 - 938^2)\,\unit{MeV^2}/(4 \times \num{6e-4}\,\unit{eV})
\approx \qty{2.7e20}{eV}$ para el fotón medio --- la cola térmica
del fondo de microondas baja el corte efectivo hasta
$\sim\qty{6e19}{eV}$. (c) Los protones por encima del corte ceden energía a
la $\Delta$ en $\sim\qty{e8}{ly}$: el espectro debería terminar cerca de
\qty{6e19}{eV} para las fuentes lejanas --- la supresión GZK observada.
(d) Sus fuentes han de ser a la vez aceleradores extraordinarios y
cósmicamente \emph{cercanas} --- dentro de nuestro supercúmulo ---, lo que
es parte de la razón por la que aún se busca su origen.
\end{solution}

\begin{solution}{exo:b3:relativistic-dynamics:12}
(a) Energía: $m_{\text{i}}c^2 = \gamma m_{\text{f}}c^2 + E_\ell$;
momento: $\gamma m_{\text{f}}\beta c = E_\ell/c$. Elimínese
$E_\ell$: $m_{\text{i}} = \gamma(1 + \beta)m_{\text{f}} =
m_{\text{f}}\sqrt{(1+\beta)/(1-\beta)}$. (b) $\sqrt{19} = 4.4$;
acelerar \emph{y} frenar lo eleva al cuadrado: $19$. (c) La masa consumida
$m_{\text{i}} - m_{\text{f}} = 3.4\,m_{\text{f}}$ debe ser
combustible aniquilado, mitad de él antimateria: \qty{1.7}{kg} de
antimateria por kilogramo entregado (solo ida, sin frenado). (d) Con una
producción mundial de nanogramos al año, los cohetes de fotones siguen
siendo aritmética, no ingeniería.
\end{solution}

\begin{solution}{pb:b3:relativistic-dynamics:1}
\textbf{1.} $E^2 - p^2c^2 = \gamma^2m^2c^4(1 - \beta^2) = m^2c^4$.
\textbf{2.} $E_{\text{tot}}$ y $\vect p_{\text{tot}}$ se transforman
como la $E$ y el $\vect p$ de una sola partícula (son sumas), de modo que el
mismo álgebra da un invariante; donde $\vect p_{\text{tot}} = \vect 0$, es
la energía total al cuadrado: $Mc^2 = E_{\text{CM}}$.
\textbf{3.} $(E + m_{\text{p}}c^2)^2 - p^2c^2 = m_{\text{p}}^2c^4 +
2Em_{\text{p}}c^2 + m_{\text{p}}^2c^4$.
\textbf{4.} $E \to m_{\text{p}}c^2$ da $M = 2m_{\text{p}}$;
$E \gg m_{\text{p}}c^2$ da $Mc^2 \to \sqrt{2Em_{\text{p}}c^2}$.
\textbf{5.} Todo movimiento relativo de los productos añade energía
cinética en el sistema del centro de momentos por encima de sus masas en
reposo; el $Mc^2$ mínimo --- el umbral --- los deja a todos en reposo
relativo.
\textbf{6.} El momento del laboratorio $\vect p_{\text{tot}} \neq \vect 0$
debe sobrevivir a la colisión: los productos están condenados a moverse, y
su energía cinética queda vedada a la fabricación de masa.
\textbf{7.} $p + p \to p + p + p + \bar p$: la carga $+2 \to +2$; el
número bariónico $+2 \to 1 + 1 + 1 - 1 = +2$. Fabricar $\bar p$ solo, o
con menos que un barión entero de más, rompe una de las dos.
\textbf{8.} Cuatro masas de protón en reposo relativo: $Mc^2 =
4m_{\text{p}}c^2$.
\textbf{9.} $16m_{\text{p}}^2c^4 = 2m_{\text{p}}^2c^4 +
2Em_{\text{p}}c^2$: $E = 7m_{\text{p}}c^2$, $T = 6m_{\text{p}}c^2 =
\qty{5.63}{GeV}$.
\textbf{10.} Masa en reposo nueva: $2m_{\text{p}}c^2$ de un $T =
6m_{\text{p}}c^2$ --- un tercio; los otros dos tercios siguen volando como
energía cinética de los productos, protegida por la conservación del momento.
\textbf{11.} Frontalmente, $Mc^2 = 2(T' + m_{\text{p}}c^2) =
4m_{\text{p}}c^2$: $T' = m_{\text{p}}c^2 = \qty{0.94}{GeV}$ por haz
--- seis veces menos que el haz sobre blanco fijo, y doce veces menos de
energía cinética total.
\textbf{12.} Los haces de los años cincuenta eran demasiado tenues para
chocar entre sí de forma útil; un blanco sólido ofrece $10^{22}$ protones
por centímetro cúbico. El precio: $\qty{5.6}{GeV}$ en lugar de
$2 \times \qty{0.94}{GeV}$.
\textbf{13.} $E = 6.2 + 0.94 = \qty{7.14}{GeV}$; $pc =
\sqrt{7.14^2 - 0.938^2} = \qty{7.08}{GeV}$.
\textbf{14.} $p = \num{7.08e9} \times \num{1.6e-19}/\num{3e8} =
\qty{3.8e-18}{kg.m/s}$: $r = p/qB = \qty{15.1}{m}$ --- la máquina tal
como se construyó.
\textbf{15.} $\beta = pc/E = 0.991$; órbita $2\pi r = \qty{95}{m}$:
$f = \qty{3.1}{MHz}$.
\textbf{16.} A \qty{10}{MeV}, $\beta = 0.145$: $f =
\qty{0.46}{MHz}$. Al subir la energía cambian tanto la velocidad (y con
ella $f$) como el momento (y con él el campo necesario a radio fijo): el
campo y la radiofrecuencia deben crecer a la par --- el truco que define al
sincrotrón.
\textbf{17.} $\num{6.19e9}/1500 \approx \num{4e6}$ vueltas; a una
frecuencia media del orden del megahercio, un ciclo de aceleración del orden
de uno o dos segundos.
\textbf{18.} El umbral del antiprotón, $T = 6m_{\text{p}}c^2 =
\qty{5.6}{GeV}$, más un margen: la contabilidad de este problema es
literalmente lo que fijó la energía --- y el nombre --- de la máquina.
\textbf{19.} $r = p/qB$ no contiene masa alguna: el imán curva por igual
momentos iguales, sea cual sea la partícula --- de modo que el haz
seleccionado sigue mezclando piones, kaones y (rara vez) antiprotones, todos
a $\qty{1.19}{GeV}/c$.
\textbf{20.} $\bar p$: $E = \sqrt{1.19^2 + 0.938^2} =
\qty{1.51}{GeV}$, $\beta = 0.79$. $\pi^-$: $E = \sqrt{1.19^2 +
0.140^2} = \qty{1.20}{GeV}$, $\beta = 0.993$.
\textbf{21.} $t_{\bar p} = 12/(0.785c) = \qty{51}{ns}$; $t_\pi =
\qty{40}{ns}$: la separación de \qty{11}{ns} exigía cronometraje de
nanosegundos --- disponible, por poco, en 1955.
\textbf{22.} Los contadores se ajustaron para dispararse con los piones
rápidos y permanecer a oscuras ante los antiprotones lentos (un veto de
velocidad); con una aguja de $1$ entre $44000$, solo dos firmas
independientes (tiempo de vuelo \emph{y} umbral Cherenkov) podían excluir
las falsas coincidencias.
\textbf{23.} Justo por encima del umbral, los productos nacen casi en
reposo relativo: la reacción apenas tiene \omterm{def:b3:hamiltonian-mechanics:hamiltonian}{espacio de fases}, mientras que la
producción de piones, muy por encima de \emph{su} umbral, es copiosa --- la
rareza quedaba garantizada por la misma cinemática que fijó el precio.
\textbf{24.} $2m_{\text{p}}c^2 = \qty{1.9}{GeV}$ por aniquilación,
típicamente en un puñado de piones que se desintegran después en fotones,
muones y neutrinos.
\textbf{25.} La conservación del \omterm{def:b3:relativistic-dynamics:four-momentum}{cuadrimomento} tasó el antiprotón en
$6m_{\text{p}}c^2 = \qty{5.6}{GeV}$ sobre un blanco fijo; el Bevatrón de
\qty{6.2}{GeV} y \qty{15.1}{m} lo pagó; el tiempo de vuelo y un veto
Cherenkov hallaron una partícula del antimundo por cada $44000$ impostores
--- y el Premio Nobel de 1959 certificó la contabilidad.
\end{solution}
